\section{Introduction}\label{sec:intro}

Ashoke Sen has published conjectures for more than thirty years, and most of them were made in the same way: a computation in one regime, a symmetry that carries it to a regime where nothing can be computed, and a statement about that second regime sharp enough to be refuted. The statement that the heterotic string on a six-torus is invariant under $SL(2,\ZZ)$ required a bound state of two monopoles that nobody had seen \cite{Sen94duality,Sen94dyon}. The statement that the tachyon of an unstable D-brane rolls to a vacuum without open strings, at an energy equal to the tension of the brane \cite{Sen98bab,Sen99univ}, required an exact solution of open string field theory that did not exist for seven years. The statement that the Fourier coefficients of a Siegel modular form have a definite sign on a cone of charges \cite{Sen11positivity} required an argument about the angular momentum of a black hole. These are predictions in the strict sense, and they can be checked.

This paper checks them. We have taken every conjecture, and every claim that was presented as more than a computation, from the papers of Sen and his collaborators that we could verify, and we have asked three questions of each: what exactly was claimed, and under which hypotheses; what has been proved since, by whom, and in which setting; and where a claim is incomplete or false, which step fails. The answers are given in the form of labelled statements. Where a conjecture has become a theorem we state the theorem and name its proof. Where it is confirmed by every test but not proved we say so, and we say what a proof would have to supply. Where it fails we give the counterexample and the hypothesis that breaks, with a computation that anyone can repeat.

The verdicts are of four kinds, and we use them as technical terms throughout.\footnote{The four verdicts are exhaustive for a statement that is precise enough to be tested. A statement that cannot be tested in any regime is recorded as \st{heuristic} in Table~\ref{tab:catalogue} and is not given a verdict; there are two of these, both about the strong coupling limit of matrix theory on tori of high dimension (Section~\ref{ssec:matrix}).}
\begin{itemize}[leftmargin=1.4em,itemsep=1pt]
\item \st{proved}: the statement, as made, follows from a published argument that we have read and that is complete within its framework. Where the framework is itself a conjecture (open string field theory as a complete theory of the open string, for instance) we say so.
\item \st{confirmed}: no proof exists and no test contradicts the statement. We name the tests, or say that there is none, and we name the missing step.
\item \st{incomplete}: the statement is correct in the regime where it was derived and has been extended, in the literature or in the statement itself, to a regime where it is not established. We give the corrected statement.
\item \st{false as stated}: there is a counterexample. We give it and we isolate the hypothesis that fails.
\end{itemize}
Two colours mark the errors in the text. \fatal{A sentence of Sen's that is fatally wrong, in the sense that the result it states falls and not merely a step of its proof, is printed in burnt orange} where we quote or restate it; the correction follows in ordinary type. \gap{A sentence that is wrong but not fatal, because it holds in the regime where it was derived and fails only in the regime to which it was extended, or because it asserts as established a step that is open, is marked with a yellow highlight.} Nothing else is coloured. Besides the two fatal errors, the yellow marks fall on the Calabi--Yau form of the Sen limit (Section~\ref{sec:senlimit}), on the Gauss--Bonnet entropy of small black holes above five dimensions (Section~\ref{ssec:small}) and on the claim that localisation derives the Rademacher expansion of the index (Section~\ref{ssec:qef}).

The result of the examination is not what a reader who starts from the suspicion that these conjectures are ``irrelevant, unproved, incomplete'' would expect, and we record it without softening it in either direction. Of the thirty-six claims of the catalogue, twenty-two are proved, ten are confirmed, two are incomplete, and two are false as stated. For every claim that is not proved we searched the literature that cites it, through October 2026, for a proof, and the verdict records what we found and the step that is missing. Statements that Sen corrected himself in a later paper are not errors of the record and are not listed. The errors are real and we treat them at length. The positivity conjecture for the index of single-centred quarter-BPS dyons \cite{Sen11positivity} fails in the freely acting $\ZZ_2$ and $\ZZ_3$ orbifolds of type II string theory on $T^6$, as Chattopadhyaya and David found \cite{CD19}; we recompute the indices from the product formula in exact arithmetic (Theorem~\ref{thm:positivityfails}), we list twenty-one charges in the single-centred cone with negative index, and we trace the failure to the explicit hypothesis of Sen's argument that a duality frame exists in which the black hole carries no hair beyond the fermion zero modes of broken supersymmetry; in these two models no such frame exists \cite{CD21}. The hypothesis holds in the heterotic string, where the conjecture was made and where it is now a theorem \cite{BringmannMurthy13,Rossello24}, and in the CHL models, where it is confirmed. The second failure is the identification of tachyon matter, the pressureless dust left at the end of the decay of an unstable brane \cite{Sen02rolling,Sen02matter,Sen02field}, as the final state of the decay. The identification is correct for homogeneous decay. For inhomogeneous decay we prove (Theorem~\ref{thm:caustic}) that a solution of the dust equations of \cite{Sen02field} develops a caustic in finite time exactly when its initial profile is not a convex function, so that every bounded profile that is not constant does, with an explicit formula for the time, and that the energy density diverges there; the effective description stops at that time, in agreement with the numerical and perturbative studies \cite{FKS02,FK04,Barnaby04}, and the closed strings emitted by the decaying brane \cite{LLM07} carry away all its energy in low dimension. What remains of the conjecture is a theorem about homogeneous decay and an open question about the inhomogeneous one, and we state both.


The two incomplete claims are the matching of the entropy of small black holes by the Gauss--Bonnet term, which holds in four and five dimensions and fails in six and more, by Sen's own computation \cite{Sen05stretching}, and the weak coupling limit of F-theory on a Calabi--Yau threefold \cite{Sen97orientifold}, whose seven-brane is singular along the orientifold plane \cite{CDE09} and whose existence fails on most bases with non-Higgsable clusters \cite{MT12,MT15}. None of these is an error in a published proof; each is a statement whose original scope was narrower than the use later made of it, and in each case the correction is in the literature, in two of the three cases in Sen's own papers.

The twenty-one proved claims include the three tachyon conjectures \cite{Schnabl05,EllwoodSchnabl07,ErlerMaccaferri14}, the existence of the two-monopole bound state \cite{Sen94dyon,Hitchin00}, the entropy function and its extension to rotating black holes \cite{Sen05entropy,AGJST06}, the wall-crossing formula for quarter-BPS dyons \cite{Sen07walls,Sen07twocentred}, which we verify exactly from the product formula (Theorem~\ref{thm:wallcrossing}), and the logarithmic corrections to the entropy of $\mathcal N=4$ and $\mathcal N=8$ black holes \cite{BGMS11}, which agree with the microscopic indices and were rederived by independent methods \cite{IMT25,IMT25b}. The ten confirmed claims include the conjecture on the $L^2$ cohomology of monopole moduli spaces, which is a theorem for two and three monopoles \cite{Hitchin00,KottkeRochon24} and partly for any number \cite{FKS18}, the logarithmic correction for $\mathcal N=2$ black holes \cite{Sen12N2}, for which no microscopic count exists, and the conjectured form of the quantum logarithmic terms in the soft graviton theorem \cite{SahooSen19}.

\paragraph*{Plan.} Section~\ref{sec:catalogue} is the catalogue, Tables~\ref{tab:catalogue} and~\ref{tab:catalogue2}, with the verdict of every claim and the place where it is treated. Sections~\ref{sec:sduality} and~\ref{sec:dyons} treat duality: the monopole conjectures and the counting of dyons, with the exact verification of the wall-crossing formula. Section~\ref{sec:positivity} treats the positivity conjecture and its failure. Sections~\ref{sec:tachyon} to~\ref{sec:nonbps} treat the tachyon conjectures, the rolling tachyon with the caustic theorem, and the non-BPS branes. Section~\ref{sec:senlimit} treats the Sen limit, Sections~\ref{sec:entropy} and~\ref{sec:quantum} the entropy function, the small black holes and the quantum entropy function, Section~\ref{sec:sft} string field theory, D-instantons, the soft theorems and matrix theory. Section~\ref{sec:conclusion} returns to the premise. Appendix~\ref{app:computation} describes the computation of the Fourier coefficients and Appendix~\ref{app:dust} completes the proof of the caustic theorem.

We have read the papers we cite. Where a journal version exists we cite it; where we could not reach a record we say so in the bibliography. We cite nothing we could not verify.\footnote{Every reference was checked against the INSPIRE or Crossref record of the paper (title, journal, volume, first page, year, DOI) during the week of 19 to 25 October 2026; the arXiv identifier is given for every paper that has one. Three references, \cite{GauntlettHarvey94}, \cite{RSZ01vacuum} and \cite{KudrnaSchnabl18}, have no journal record in INSPIRE and are cited as preprints.}
